\documentclass[11pt,a4paper]{article}

\usepackage[utf8]{inputenc}
\usepackage[T1]{fontenc}
\usepackage[french]{babel}
\usepackage[a4paper,margin=2.2cm]{geometry}
\usepackage{amsmath,amssymb,amsthm}
\usepackage{lmodern}
\usepackage{enumitem}
\usepackage{hyperref}
\usepackage{fancyhdr}
\usepackage{mathrsfs}

\hypersetup{
    colorlinks=true,
    linkcolor=black,
    urlcolor=blue
}

\setlength{\parindent}{0pt}
\setlength{\parskip}{0.5em}

\pagestyle{fancy}
\fancyhf{}
\fancyhead[L]{Devoir surveillé de mathématiques-PGB}
\fancyhead[R]{Terminale --- Analyse fonctionnelle}
\fancyfoot[C]{\thepage}

\newtheorem*{remarque}{Remarque}

\newcommand{\R}{\mathbb{R}}
\newcommand{\C}{\mathbb{C}}
\newcommand{\cC}{\mathcal{C}}

\begin{document}

\begin{center}
    {\LARGE \textbf{Devoir surveillé de mathématiques}}\\[0.5em]
    {\large \textbf{Niveau : Terminale générale / Maths expertes possible}}\\[0.5em]
    \textbf{Durée conseillée : 4 heures}
\end{center}

\hrule
\vspace{0.8em}

\textbf{Consignes générales}

La qualité de la rédaction, la précision des justifications et la clarté du raisonnement seront prises en compte de façon importante dans la notation.

Toujours pas de calculatrice, Pas de CHATGPT ni aucune aide extérieure.
IMPORTANTTTTTTTTT 
Pensez à noter en haut de votre 1ere copie le temps que vous avez pris à faire le DM.

Aucune initiative pertinente ne sera pénalisée, même si elle n’aboutit pas complètement.

\vspace{1em}
\hrule
\vspace{1em}

\section*{Problème 1 --- Étude d’une fonction définie par ses propriétés locales}

Dans tout le problème, on considère une fonction \(f\) définie sur \(\R\), continue sur \(\R\), dérivable sur \(\R\), vérifiant les propriétés suivantes :
\begin{enumerate}[label=\arabic*.]
    \item \(f(0)=0\),
    \item pour tous réels \(x\) et \(h\) avec \(h\neq 0\),
    \[
    \frac{f(x+h)-f(x)}{h} \ge 1-\frac{h^2}{1+h^2},
    \]
    \item pour tout réel \(x\),
    \[
    f(-x)=-f(x).
    \]
\end{enumerate}



\subsection*{Partie A --- Premiers résultats}

\begin{enumerate}[label=\arabic*.]
    \item Déterminer la parité de \(f'\), lorsqu’elle existe.
    \item Montrer que, pour tout réel \(x\), on a \(f'(x)\ge 1\).
    \item En déduire les variations de \(f\) sur \(\R\).
    \item Montrer que, pour tout \(x\ge 0\), on a \(f(x)\ge x\).
    \item Que peut-on en déduire pour \(x\le 0\) ?
\end{enumerate}



\subsection*{Partie B --- Encadrement, convexité intuitive, comparaison de taux de variation}

On suppose dans cette partie que, pour tous réels \(x<y<z\),
\[
\frac{f(y)-f(x)}{y-x}\le \frac{f(z)-f(y)}{z-y}.
\]

\begin{enumerate}[label=\arabic*.]
    \item Montrer que pour tous réels \(a<b\),
    \[
    f\!\left(\frac{a+b}{2}\right)\le \frac{f(a)+f(b)}{2}.
    \]
    \item Montrer que pour tout réel \(x\),
    \[
    f(x)+f(-x)\ge 2f(0).
    \]
    \item Confronter ce résultat à la propriété de parité de \(f\).
    \item Que peut-on conclure sur \(f(0)\) ? Ce résultat était-il déjà connu ?
    
\end{enumerate}

\subsection*{Partie D --- Comportement aux bornes et limites(la plus facile)}

On conserve toutes les hypothèses des parties précédentes.

\begin{enumerate}[label=\arabic*.]
    

    \item Montrer que
    \[
    \lim_{x\to +\infty} f(x)=+\infty.
    \]

    

    \item Montrer aussi du coup que
    \[
    \lim_{x\to -\infty} f(x)=-\infty.
    \]

     On pose
    \[
    \varphi(x)=f(x)-x.
    \]
    

   

    \item Étudier la parité de \(\varphi\).

   
\end{enumerate}

\vspace{1em}
\hrule
\vspace{1em}

\section*{Problème 2 --- Décomposition paire/impaire et structure de l’ensemble des fonctions continues}

On note \(\cC(\R)\) l’ensemble des fonctions continues de \(\R\) dans \(\R\).

Pour toute fonction \(f\in\cC(\R)\), on définit les fonctions \(p_f\) et \(i_f\) par
\[
p_f(x)=\frac{f(x)+f(-x)}{2},
\qquad
i_f(x)=\frac{f(x)-f(-x)}{2}.
\]

\subsection*{Partie A --- Mise en place}

\begin{enumerate}[label=\arabic*.]
    \item Montrer que \(p_f\) et \(i_f\) appartiennent à \(\cC(\R)\).
    \item Montrer que \(p_f\) est paire et que \(i_f\) est impaire.
    \item Montrer que
    \[
    f=p_f+i_f.
    \]
\end{enumerate}

\subsection*{Partie B --- Unicité de la décomposition}

\begin{enumerate}[label=\arabic*.]
    \item Soient \(u\) une fonction paire et \(v\) une fonction impaire telles que
    \[
    f=u+v.
    \]
    Montrer que \(u=p_f\) et \(v=i_f\).
    \item En déduire que toute fonction continue sur \(\R\) s’écrit d’une et une seule manière comme somme d’une fonction paire et d’une fonction impaire.
    \item Montrer qu’une fonction à la fois paire et impaire est nécessairement nulle.
    \item Expliquer en quoi ce résultat est lié à l’unicité précédente.
\end{enumerate}

\subsection*{Partie C --- Applications analytiques}

On considère désormais une fonction \(f\) continue sur \(\R\).

\begin{enumerate}[label=\arabic*.]
    \item Montrer que si \(f\) est dérivable en \(0\) et paire, alors \(f'(0)=0\).
    \item Montrer que si \(f\) est dérivable et impaire, alors \(f'\) est paire.
    \item Montrer que si \(f\) est dérivable et paire, alors \(f'\) est impaire.
    \item Soit \(f\) dérivable sur \(\R\).  
    Exprimer \(p_{f'}\) et \(i_{f'}\) à l’aide de \(p_f\) et \(i_f\).
    \item On suppose maintenant \(f\) deux fois dérivable.  
    Discuter la parité de \(f''\) selon celle de \(f\).
\end{enumerate}

\subsection*{Partie D --- c bientot finis: application directe}

On considère la fonction \(g\) définie sur \(\R\) par
\[
g(x)=x^3-3x^2+4x.
\]

\begin{enumerate}[label=\arabic*.]
    \item Calculer \(g(-x)\).  
    La fonction \(g\) est-elle paire, impaire, ou ni paire ni impaire ?

    \item Déterminer la partie paire \(p_g\) et la partie impaire \(i_g\) de \(g\).

    \item Vérifier que
    \[
    g=p_g+i_g.
    \]

    \item Montrer que \(p_g\) est paire et que \(i_g\) est impaire.

    \item Calculer \(g'(x)\).

    \item Étudier le signe de \(g'(x)\).

    \item En déduire les variations de \(g\) sur \(\R\).

    \item Dresser le tableau de variation de \(g\).

    \item Déterminer les extremums éventuels de \(g\).

    \item Déterminer les limites de \(g\) en \(+\infty\) et en \(-\infty\).

    \item Calculer \(p_g'(x)\) et \(i_g'(x)\).

    \item Étudier les variations de \(p_g\) sur \([0,+\infty[\).

    \item Étudier les variations de \(i_g\) sur \(\R\).

    \item Comparer l’allure des courbes représentatives de \(g\), \(p_g\) et \(i_g\).
\end{enumerate}
\subsection*{Partie E --- Étude d’une condition fonctionnelle}

On cherche les fonctions continues \(f\) sur \(\R\) telles que, pour tout réel \(x\),
\[
f(x)+f(-x)=2x^2.
\]

\begin{enumerate}[label=\arabic*.]
    \item Montrer qu’une telle fonction ne peut pas être impaire.
    \item Déterminer sa partie paire.
    \item Montrer que l’ensemble des solutions peut s’écrire à l’aide d’une fonction impaire arbitraire.
    \item Donner trois exemples distincts de telles fonctions.
    \item Parmi elles, déterminer celles qui sont dérivables sur \(\R\).
    \item Parmi elles, déterminer celles qui sont monotones sur \([0,+\infty[\).
\end{enumerate}

\vspace{1em}
\hrule
\vspace{1em}



\section*{Exercice Bonus --- Taux de variation sans expression explicite}

Soit \(f\) une fonction définie sur \(\R\), et soit \(a\in\R\).
\varepsilon\) est une fonction réelle
On suppose que, pour tout réel \(h\neq 0\),
\[
\frac{f(a+h)-f(a)}{h}=2a+h+\varepsilon(h),
\]
où \(\varepsilon(h)\to 0\) quand \(h\to 0\).

\begin{enumerate}[label=\arabic*.]
    \item Montrer que \(f\) est dérivable en \(a\).
    \item Déterminer \(f'(a)\).
    \item Interpréter précisément le rôle du terme \(2a+h+\varepsilon(h)\).
    \item Cette information permet-elle de connaître l’expression de \(f\) ? Justifier soigneusement.
\end{enumerate}

\end{document}